\documentclass[11pt,a4paper]{article}
\usepackage[utf8]{inputenc}
\usepackage[T1]{fontenc}
\usepackage[margin=2.3cm]{geometry}
\usepackage{graphicx}
\usepackage{booktabs}
\usepackage{xcolor}
\usepackage{hyperref}
\definecolor{copper}{HTML}{C96F4A}
\hypersetup{colorlinks=true, linkcolor=copper, urlcolor=copper}
\setlength{\parskip}{4pt}
\begin{document}
\begin{center}
{\LARGE\bfseries Research Program: The Machine with No Long-Term Memory (ES1–ES5, TS1–TS3)}\\[8pt]
{\large Research Program Monograph · Version 1.0.0}\\[6pt]
{\color{copper}\small RESEARCH PROGRAM · SPINOR}\\[4pt]
{\small Isaev Iskhak Khamzatovich · ORCID 0009-0003-7299-0701}\\[2pt]
{\small AB Cloud Quantum Research · DOI 10.5281/zenodo.21825394}
\end{center}

\tableofcontents
\newpage

\section{Research Program: The Spinor Study}

\noindent\textit{\color{copper}The spinor program drives the central claim "reads, does not carry" to explicit models: the only state is the spinor phase; a step is a rotation by θ. Three theorems, five campaigns, all PASS; the capacity-versus-robustness trade-off is fixed as an honest boundary.}

\subsection{1. Program goals}

\begin{itemize}
  \item state and prove the properties of the volatile state: cycle closure (TS1), reset irreversibility (TS2), O(1) steps versus the Ω(t) log (TS3);
  \item measure the models with campaigns carrying pre-registered thresholds;
  \item prepare the models for processor stage P1 (memory V21).
\end{itemize}

\subsection{2. Models}

\begin{table}[htbp]\centering\small
\caption{Table 4.2. Three models and their parameters}
\begin{tabular}{p{0.10\textwidth} p{0.10\textwidth} p{0.13\textwidth} p{0.15\textwidth} p{0.13\textwidth} p{0.35\textwidth}}
\toprule
\textbf{Model} & \textbf{θ} & \textbf{N = 2k} & \textbf{Bits/step} & \textbf{σ*} & \textbf{Purpose}
\midrule
S1 & π/2 & 4 & 2.000 & 0.135 & maximal robustness \\
S2 & π/3 & 6 & 2.585 & 0.089 & balance \\
S3 & π/7 & 14 & 3.807 & 0.038 & maximal capacity \\
\bottomrule
\end{tabular}
\end{table}

\subsection{3. Executed campaigns ES1–ES5}

\begin{table}[htbp]\centering\small
\caption{Table 4.3. Spinor campaigns}
\begin{tabular}{p{0.12\textwidth} p{0.19\textwidth} p{0.52\textwidth} p{0.13\textwidth}}
\toprule
\textbf{Campaign} & \textbf{Measured} & \textbf{Result} & \textbf{Verdict}
\midrule
ES1 & cycles & periods 4/6/14 = 2k exactly; gap = step & PASS \\
ES2 & reset & I(pre;post) = 0; reconstruction ≈ 1/N & PASS \\
ES3 & time & step slope −1.5e−12; log ×129 & PASS \\
ES4 & task & a Markov policy is solvable without memory & PASS \\
ES5 & noise & diffusion bounded by the silence zone & PASS \\
\bottomrule
\end{tabular}
\end{table}

The errata ES2a and ES5a: the thresholds of campaigns ES2 and ES5 were replaced before the verdicts were accepted; the replacements are published in the verdict files. The log₂N versus 1/k trade-off is a permanent boundary: capacity grows, the silence zone narrows.

\subsection{4. The next wave}

\begin{itemize}
  \item phases θ = π/k for k > 7: searching for the inflection points of the trade-off;
  \item task compilation into phase cycles: a formal grammar of policies;
  \item a diffusion channel with correlations: beyond the Markov assumption.
\end{itemize}

\subsection{5. Honest boundaries}

Not every task is solvable without memory: the applicability class is tasks with a Markov policy. Reconstruction ≈ 1/N is the chance level, not "almost zero": information is erased to the limit. A log gives Ω(t) under any "smart" scheduling — this is a lower bound, not an implementation defect.

\subsection{6. Campaigns in detail}

\subsection{6.1. Five measurements of volatility}

\subsubsection{ES1 — cycles}

The first measurement campaign: the rotation periods for θ = π/2, π/3, π/7 equal 4, 6 and 14 — exactly N = 2k, and the gap between distinguishable states equals the step 2π/N. The volatile phase turned out to be an exact counter: no drift, no memory.

\subsubsection{ES2 — reset and erasure}

A phase reset erases everything: the pre/post mutual information is zero, reconstruction 0.2501/0.1674/0.0717 ≈ 1/N — the chance level. Time after the reset is defined only modulo N. Theorem TS2 gets its numerical body here; erratum ES2a recorded the threshold replacement before the verdict.

\subsubsection{ES3 — real time}

The contrast experiment: the memoryless machine holds O(1) steps (slope −1.5e−12 ms/step — zero), the eternal-log machine grows ×129 and misses deadlines. This is the numerical proof of TS3 and the main argument of the "reads, does not carry" architecture: a log is not reliability but a cause of failure.

\subsubsection{ES4 — the memoryless task}

The applicability check: a sequence of decisions inside the cycle without storage; on the class of tasks with a Markov policy the quality is comparable to a memory-bearing machine. The boundary is honestly fixed: tasks with substantial history require another level of the architecture.

\subsubsection{ES5 — noise and the silence zone}

Phase diffusion is bounded by the silence zone σ*: S1 (0.135 rad) is more robust, S3 (0.038 rad) has more capacity — the capacity/robustness trade-off is measured. Erratum ES5a: the threshold was replaced before the verdict was accepted, and the replacement is published.

\subsection{7. The full campaign table}

\begin{table}[htbp]\centering\small
\caption{The full campaign table of the level}
\begin{tabular}{p{0.12\textwidth} p{0.25\textwidth} p{0.60\textwidth}}
\toprule
\textbf{Campaign} & \textbf{Title} & \textbf{Key result}
\midrule
ES1 & Phase cycles &  \\
ES2 & Volatility &  \\
ES3 & Real time &  \\
ES4 & Memoryless task &  \\
ES5 & Noise diffusion &  \\
\bottomrule
\end{tabular}
\end{table}

\subsection{8. How to reproduce}

\begin{verbatim}
cd research/realtime_spinor_memory
python3 experiments/run_all.py        # ES1–ES5 (~2.5 min, all PASS)
python3 make_figures.py               # 4-dimension tiles
python3 make_theorem_docs.py          # 12 TS1–TS3 documents
python3 -m pytest tests/ -q           # 11 tests
\end{verbatim}

\subsection{9. Further reading}

\begin{itemize}
  \item theory: docs/theory.en.md (EN) and docs/theory.md (RU);
  \item implementation: spinor/core.py — the model core, spinor/tasks.py — tasks;
  \item verdicts: data/es1…es5\_verdict.json;
  \item theorems: theorems/TS1\_cycle, TS2\_volatility, TS3\_realtime.
\end{itemize}

\end{document}
